% Nouveau EGA V (cote 25): archive cote 25, site page 156 (PDF page 157).
% typescript. Draft compared: none.
% AI-transcribed by Claude Opus 5.5 (claude-opus-5-5, adaptive thinking, effort high), Message Batches API, batch ega25b (msgbatch_01A7n5YuXubUPNK1zsepi8RG), 2026-11-13; not yet checked by eye.
\scanpage{25/156}
\mspage{\hand{17}}

\xout{xg} ``générique'')\hand{\unsure{d}}.   Ceci posé, il faudrait vérifier que l'ensem ble $V^2$
des points supersinguliers de u rel. à Y n'est autre que l'ensemble des
points en lesquels $V^{(1)} \longrightarrow Y$ est ramifié, et préciser ce point
comme une identité de sous-préschémas fermés de X , $V^{\unsure{\over{2}{2}}}$ \over{=}{=} \hand{``}sous-préschéma
de ramification\hand{''} de $V^1 \longrightarrow Y$. Je n'ai pas fait l'exercice en détail,
mais ne doute pas qu'il ne puisse se faire par essentiellement \over{la}{le} même
\xout{méthode} \inserted{genre de dévissage} que cel\over{le}{ui} développé\deleted{e} après la prop.7.

\hand{\unsure{d}} Si $E \longrightarrow F$ est un homomorphisme de Modules localement libres \unsure{sur} \inserted{le préschéma régulier}
X, de rangs m et d ($m \over{<}{\leqslant} d$), et si \over{\unsure{Y}}{Z} est l'ensemble des points de X en
lesquels le rang de u est $\hand{\leqslant}$ m-1, alors \inserted{on peut montrer que} Z est de codimension $\hand{\geqslant}$ d-m+1
en chacun de ses points, \xout{comme on vérifie immédiatement par récurrence sur}
\xout{m, à partir du cas m=1 (}

\mspage{\hand{156}}
